\chapter{Legal Contracts: Adjudication as Non-Informational Resolution}
\label{ch:contracts}

\begin{quotation}
\noindent\textit{Synopsis.} Contracts attempt to regulate future
contingencies through finite language, making them a natural example of
continuation under uncertainty. This chapter examines how legal systems
respond when a contract encounters a situation its authors did not
anticipate, and presents adjudication as a distinctive mechanism: a court
resolves ambiguity through authoritative decision rather than by
recovering additional distinctions about the past. Law is, on this
framework, a domain in which continuation proceeds despite persistent,
unresolved informational insufficiency.
\end{quotation}

\section{Contracts as Attempted Brute-Force Closure}

A contract's enumerated contingencies (``in the event of $X$, $Y$, or
$Z$'') attempt to close what is, in fact, an open contingency space by
brute enumeration, resembling augmentation (\cref{sec:coord-augmentation})
in strategy while targeting a domain that does not satisfy
\cref{def:closed-open}'s closed case.

\begin{example}
\label{ex:force-majeure}
A force majeure clause enumerating war, natural disaster, and government
action as excusing conditions is an augmentation set $H_{\text{contract}}
= \{$war, natural disaster, government action$\}$ in exactly the sense of
\cref{prop:closed-sufficiency}: it pre-commits the parties' agreement to
cover these named contingencies without needing to anticipate, at the
time of drafting, which one might occur. The clause is sufficient for
any dispute falling within $H_{\text{contract}}$. A disruption outside
this list -- one neither party had reason to name at drafting time --
is not covered, not because the clause was drafted carelessly, but
because $H_{\text{contract}}$, like any finite augmentation, cannot
exhaust an open space of future circumstances.
\end{example}

\begin{proposition}[Contractual Augmentation Targets the Wrong Kind of Domain]
\label{prop:contract-augmentation-mismatch}
A contract's enumerated-contingency strategy achieves
\cref{prop:closed-sufficiency}'s sufficiency only for contingencies
within its enumerated set $H_{\text{contract}}$; by
\cref{prop:sufficiency-closed-only}, no finite $H_{\text{contract}}$
achieves sufficiency for the full space of future circumstances, since
that space is open in the sense of \cref{def:closed-open}.
\end{proposition}

\begin{proofsketch}
Immediate from \cref{prop:sufficiency-closed-only} applied with the
future-circumstance space playing the role of the open hidden-state
space: augmentation is the correct strategy in form, exactly as it is
for chess (\cref{ch:chess}), but the domain it is applied to lacks the
closure property chess's hidden state has, so the strategy's guarantee
does not transfer.
\end{proofsketch}

This mismatch between strategy and domain is, on the framework developed
here, a structural explanation for the persistence of contractual
litigation: no finite enumeration of contingencies exhausts an open
space of future circumstances, and every contract that enumerates
contingencies is, in this precise sense, applying chess's solution to
painting's problem.

\section{Adjudication, Worked}

\cref{prop:adjudication} is developed here in full: a court's ruling on a
disputed clause selects a binding continuation without recovering any
additional distinguishability information about the parties' original
intent. The information deficit -- what the parties actually meant at
the time of drafting -- remains after the ruling exactly as before it;
what has changed is that further continuation is now permitted to
proceed on the ruling's basis. In the vocabulary of
Section~\ref{sec:coord-adjudication}, the court substitutes its ruling as
the operative correctness set $\mathcal{C}_\tau(e)$ going forward,
regardless of whether that ruling coincides with whatever the
undiscoverable true correctness set would have contained.

\section{Litigating Intent as Reconstitution}

Establishing original intent is a reconstitution-mode task
(\cref{ch:mode-relativity}): for a court asking about intent, $r_{c,t}$
is correspondence, drafting history, and testimony, none of which
constitutes direct access to the intent itself
(\cref{def:available-substrate}). This is a further instance of
\cref{prop:recursive-insufficiency}: the traces used to reconstruct
intent are themselves finite substrates subject to
\cref{thm:composition-failure}, which is part of why intent disputes
resist final resolution by evidence alone and are ultimately settled,
per \cref{prop:adjudication}, by authority rather than by information
recovery.
